\chapter{\'Evaluation}
\label{ch:evaluation}

\section{Protocole exp\'erimental}

Les exp\'eriences ont \'et\'e men\'ees sur un cluster de cinq
n\oe{}uds. Le tableau~\ref{tab:resultats} r\'esume le d\'ebit obtenu
pour des taux de partition croissants.

\begin{table}[htbp]
  \centering
  \caption{D\'ebit (requ\^etes/s) selon le taux de partition simul\'e.}
  \label{tab:resultats}
  \begin{tabular}{@{}lrr@{}}
    \toprule
    Taux de partition & R\'ef\'erence & Propos\'e \\
    \midrule
    0\%  & 12\,400 & 12\,100 \\
    5\%  & 9\,800  & 11\,600 \\
    10\% & 6\,200  & 10\,900 \\
    \bottomrule
  \end{tabular}
\end{table}

\section{Latence sous charge}

La figure~\ref{fig:latence} montre la latence m\'ediane des requ\^etes
lorsque le taux de partition augmente.

\begin{figure}[htbp]
  \centering
  \rule{0.6\textwidth}{4cm}
  \caption{Latence m\'ediane en fonction du taux de partition simul\'e.}
  \label{fig:latence}
\end{figure}

\section{Discussion}

Le protocole propos\'e maintient un meilleur d\'ebit que la
r\'ef\'erence d\`es que le taux de partition d\'epasse 5~\%, au prix
d'un faible surco\^ut de latence constant \citep{ongaro2014search}.
